\documentclass[11pt,a4paper]{article}
% lesson-preamble.tex -- Shared preamble for all Phase 00 lessons
% Usage: % lesson-preamble.tex -- Shared preamble for all Phase 00 lessons
% Usage: % lesson-preamble.tex -- Shared preamble for all Phase 00 lessons
% Usage: \input{../lesson-preamble} (from subdomain directories)

% ---- Encoding and fonts ----
\usepackage[utf8]{inputenc}
\usepackage[T1]{fontenc}

% ---- Mathematics ----
\usepackage{amsmath,amssymb,amsthm}
\usepackage{mathtools}

% ---- Page geometry ----
\usepackage[margin=1in,headheight=14pt]{geometry}

% ---- Hyperlinks ----
\usepackage[colorlinks=true,linkcolor=red,citecolor=blue,urlcolor=blue]{hyperref}

% ---- Lists ----
\usepackage{enumitem}

% ---- Colors and boxes ----
\usepackage{xcolor}
\usepackage[theorems,skins,breakable]{tcolorbox}

% ---- Header/Footer ----
\usepackage{fancyhdr}
\renewcommand{\headrulewidth}{0.4pt}

% ---- Tables ----
\usepackage{booktabs}

% ---- Bibliography ----
\usepackage{natbib}
\bibliographystyle{plainnat}

% --------------------------------------------------------------------------
% Theorem environments
% --------------------------------------------------------------------------
\theoremstyle{plain}
\newtheorem{theorem}{Theorem}[section]
\newtheorem{lemma}[theorem]{Lemma}
\newtheorem{proposition}[theorem]{Proposition}
\newtheorem{corollary}[theorem]{Corollary}

\theoremstyle{definition}
\newtheorem{definition}[theorem]{Definition}
\newtheorem{example}[theorem]{Example}
\newtheorem{exercise}[theorem]{Exercise}
\newtheorem{assumption}[theorem]{Assumption}

\theoremstyle{remark}
\newtheorem{remark}[theorem]{Remark}

% --------------------------------------------------------------------------
% Colored boxes
% --------------------------------------------------------------------------
\newtcolorbox{keyresult}[1][]{%
  colback=green!5!white,
  colframe=green!50!black,
  fonttitle=\bfseries,
  title={Key Result},
  breakable,
  #1
}

\newtcolorbox{rlconnection}[1][]{%
  colback=blue!5!white,
  colframe=blue!50!black,
  fonttitle=\bfseries,
  title={RL/ML Connection},
  breakable,
  #1
}

\newtcolorbox{warning}[1][]{%
  colback=red!5!white,
  colframe=red!50!black,
  fonttitle=\bfseries,
  title={Warning},
  breakable,
  #1
}

% --------------------------------------------------------------------------
% Custom commands -- number sets
% --------------------------------------------------------------------------
\newcommand{\R}{\mathbb{R}}
\newcommand{\N}{\mathbb{N}}
\newcommand{\Q}{\mathbb{Q}}
\newcommand{\Z}{\mathbb{Z}}
\newcommand{\C}{\mathbb{C}}
\newcommand{\F}{\mathbb{F}}

% --------------------------------------------------------------------------
% Custom commands -- probability and expectation
% --------------------------------------------------------------------------
\newcommand{\E}{\mathbb{E}}
\newcommand{\Prob}{\mathbb{P}}
\newcommand{\Var}{\operatorname{Var}}
\newcommand{\Cov}{\operatorname{Cov}}
\newcommand{\indicator}{\mathbf{1}}

% --------------------------------------------------------------------------
% Custom commands -- convergence arrows
% --------------------------------------------------------------------------
\newcommand{\cas}{\xrightarrow{\text{a.s.}}}
\newcommand{\cprob}{\xrightarrow{\Prob}}
\newcommand{\clp}[1]{\xrightarrow{L^{#1}}}
\newcommand{\cdist}{\xrightarrow{d}}

% --------------------------------------------------------------------------
% Custom commands -- common operators
% --------------------------------------------------------------------------
\DeclareMathOperator{\dom}{dom}
\DeclareMathOperator{\epi}{epi}
\DeclareMathOperator{\conv}{conv}
\DeclareMathOperator{\relint}{relint}
\DeclareMathOperator{\aff}{aff}
\DeclareMathOperator{\cl}{cl}
\DeclareMathOperator{\interior}{int}
\DeclareMathOperator{\tr}{tr}
\DeclareMathOperator{\diag}{diag}
\DeclareMathOperator{\rank}{rank}
\DeclareMathOperator{\sgn}{sgn}
\DeclareMathOperator{\supp}{supp}
\DeclareMathOperator{\argmin}{arg\,min}
\DeclareMathOperator{\argmax}{arg\,max}
\DeclareMathOperator{\prox}{prox}
\DeclareMathOperator{\dist}{dist}
\DeclareMathOperator{\diam}{diam}
\DeclareMathOperator{\Span}{span}

% --------------------------------------------------------------------------
% Custom commands -- linear algebra operators
% --------------------------------------------------------------------------
\DeclareMathOperator{\nullity}{nullity}
\DeclareMathOperator{\range}{range}
\DeclareMathOperator{\nullspace}{null}
\DeclareMathOperator{\proj}{proj}

% --------------------------------------------------------------------------
% Custom commands -- measure theory
% --------------------------------------------------------------------------
\newcommand{\powerset}{\mathcal{P}}
\newcommand{\Borel}{\mathcal{B}}
 (from subdomain directories)

% ---- Encoding and fonts ----
\usepackage[utf8]{inputenc}
\usepackage[T1]{fontenc}

% ---- Mathematics ----
\usepackage{amsmath,amssymb,amsthm}
\usepackage{mathtools}

% ---- Page geometry ----
\usepackage[margin=1in,headheight=14pt]{geometry}

% ---- Hyperlinks ----
\usepackage[colorlinks=true,linkcolor=red,citecolor=blue,urlcolor=blue]{hyperref}

% ---- Lists ----
\usepackage{enumitem}

% ---- Colors and boxes ----
\usepackage{xcolor}
\usepackage[theorems,skins,breakable]{tcolorbox}

% ---- Header/Footer ----
\usepackage{fancyhdr}
\renewcommand{\headrulewidth}{0.4pt}

% ---- Tables ----
\usepackage{booktabs}

% ---- Bibliography ----
\usepackage{natbib}
\bibliographystyle{plainnat}

% --------------------------------------------------------------------------
% Theorem environments
% --------------------------------------------------------------------------
\theoremstyle{plain}
\newtheorem{theorem}{Theorem}[section]
\newtheorem{lemma}[theorem]{Lemma}
\newtheorem{proposition}[theorem]{Proposition}
\newtheorem{corollary}[theorem]{Corollary}

\theoremstyle{definition}
\newtheorem{definition}[theorem]{Definition}
\newtheorem{example}[theorem]{Example}
\newtheorem{exercise}[theorem]{Exercise}
\newtheorem{assumption}[theorem]{Assumption}

\theoremstyle{remark}
\newtheorem{remark}[theorem]{Remark}

% --------------------------------------------------------------------------
% Colored boxes
% --------------------------------------------------------------------------
\newtcolorbox{keyresult}[1][]{%
  colback=green!5!white,
  colframe=green!50!black,
  fonttitle=\bfseries,
  title={Key Result},
  breakable,
  #1
}

\newtcolorbox{rlconnection}[1][]{%
  colback=blue!5!white,
  colframe=blue!50!black,
  fonttitle=\bfseries,
  title={RL/ML Connection},
  breakable,
  #1
}

\newtcolorbox{warning}[1][]{%
  colback=red!5!white,
  colframe=red!50!black,
  fonttitle=\bfseries,
  title={Warning},
  breakable,
  #1
}

% --------------------------------------------------------------------------
% Custom commands -- number sets
% --------------------------------------------------------------------------
\newcommand{\R}{\mathbb{R}}
\newcommand{\N}{\mathbb{N}}
\newcommand{\Q}{\mathbb{Q}}
\newcommand{\Z}{\mathbb{Z}}
\newcommand{\C}{\mathbb{C}}
\newcommand{\F}{\mathbb{F}}

% --------------------------------------------------------------------------
% Custom commands -- probability and expectation
% --------------------------------------------------------------------------
\newcommand{\E}{\mathbb{E}}
\newcommand{\Prob}{\mathbb{P}}
\newcommand{\Var}{\operatorname{Var}}
\newcommand{\Cov}{\operatorname{Cov}}
\newcommand{\indicator}{\mathbf{1}}

% --------------------------------------------------------------------------
% Custom commands -- convergence arrows
% --------------------------------------------------------------------------
\newcommand{\cas}{\xrightarrow{\text{a.s.}}}
\newcommand{\cprob}{\xrightarrow{\Prob}}
\newcommand{\clp}[1]{\xrightarrow{L^{#1}}}
\newcommand{\cdist}{\xrightarrow{d}}

% --------------------------------------------------------------------------
% Custom commands -- common operators
% --------------------------------------------------------------------------
\DeclareMathOperator{\dom}{dom}
\DeclareMathOperator{\epi}{epi}
\DeclareMathOperator{\conv}{conv}
\DeclareMathOperator{\relint}{relint}
\DeclareMathOperator{\aff}{aff}
\DeclareMathOperator{\cl}{cl}
\DeclareMathOperator{\interior}{int}
\DeclareMathOperator{\tr}{tr}
\DeclareMathOperator{\diag}{diag}
\DeclareMathOperator{\rank}{rank}
\DeclareMathOperator{\sgn}{sgn}
\DeclareMathOperator{\supp}{supp}
\DeclareMathOperator{\argmin}{arg\,min}
\DeclareMathOperator{\argmax}{arg\,max}
\DeclareMathOperator{\prox}{prox}
\DeclareMathOperator{\dist}{dist}
\DeclareMathOperator{\diam}{diam}
\DeclareMathOperator{\Span}{span}

% --------------------------------------------------------------------------
% Custom commands -- linear algebra operators
% --------------------------------------------------------------------------
\DeclareMathOperator{\nullity}{nullity}
\DeclareMathOperator{\range}{range}
\DeclareMathOperator{\nullspace}{null}
\DeclareMathOperator{\proj}{proj}

% --------------------------------------------------------------------------
% Custom commands -- measure theory
% --------------------------------------------------------------------------
\newcommand{\powerset}{\mathcal{P}}
\newcommand{\Borel}{\mathcal{B}}
 (from subdomain directories)

% ---- Encoding and fonts ----
\usepackage[utf8]{inputenc}
\usepackage[T1]{fontenc}

% ---- Mathematics ----
\usepackage{amsmath,amssymb,amsthm}
\usepackage{mathtools}

% ---- Page geometry ----
\usepackage[margin=1in,headheight=14pt]{geometry}

% ---- Hyperlinks ----
\usepackage[colorlinks=true,linkcolor=red,citecolor=blue,urlcolor=blue]{hyperref}

% ---- Lists ----
\usepackage{enumitem}

% ---- Colors and boxes ----
\usepackage{xcolor}
\usepackage[theorems,skins,breakable]{tcolorbox}

% ---- Header/Footer ----
\usepackage{fancyhdr}
\renewcommand{\headrulewidth}{0.4pt}

% ---- Tables ----
\usepackage{booktabs}

% ---- Bibliography ----
\usepackage{natbib}
\bibliographystyle{plainnat}

% --------------------------------------------------------------------------
% Theorem environments
% --------------------------------------------------------------------------
\theoremstyle{plain}
\newtheorem{theorem}{Theorem}[section]
\newtheorem{lemma}[theorem]{Lemma}
\newtheorem{proposition}[theorem]{Proposition}
\newtheorem{corollary}[theorem]{Corollary}

\theoremstyle{definition}
\newtheorem{definition}[theorem]{Definition}
\newtheorem{example}[theorem]{Example}
\newtheorem{exercise}[theorem]{Exercise}
\newtheorem{assumption}[theorem]{Assumption}

\theoremstyle{remark}
\newtheorem{remark}[theorem]{Remark}

% --------------------------------------------------------------------------
% Colored boxes
% --------------------------------------------------------------------------
\newtcolorbox{keyresult}[1][]{%
  colback=green!5!white,
  colframe=green!50!black,
  fonttitle=\bfseries,
  title={Key Result},
  breakable,
  #1
}

\newtcolorbox{rlconnection}[1][]{%
  colback=blue!5!white,
  colframe=blue!50!black,
  fonttitle=\bfseries,
  title={RL/ML Connection},
  breakable,
  #1
}

\newtcolorbox{warning}[1][]{%
  colback=red!5!white,
  colframe=red!50!black,
  fonttitle=\bfseries,
  title={Warning},
  breakable,
  #1
}

% --------------------------------------------------------------------------
% Custom commands -- number sets
% --------------------------------------------------------------------------
\newcommand{\R}{\mathbb{R}}
\newcommand{\N}{\mathbb{N}}
\newcommand{\Q}{\mathbb{Q}}
\newcommand{\Z}{\mathbb{Z}}
\newcommand{\C}{\mathbb{C}}
\newcommand{\F}{\mathbb{F}}

% --------------------------------------------------------------------------
% Custom commands -- probability and expectation
% --------------------------------------------------------------------------
\newcommand{\E}{\mathbb{E}}
\newcommand{\Prob}{\mathbb{P}}
\newcommand{\Var}{\operatorname{Var}}
\newcommand{\Cov}{\operatorname{Cov}}
\newcommand{\indicator}{\mathbf{1}}

% --------------------------------------------------------------------------
% Custom commands -- convergence arrows
% --------------------------------------------------------------------------
\newcommand{\cas}{\xrightarrow{\text{a.s.}}}
\newcommand{\cprob}{\xrightarrow{\Prob}}
\newcommand{\clp}[1]{\xrightarrow{L^{#1}}}
\newcommand{\cdist}{\xrightarrow{d}}

% --------------------------------------------------------------------------
% Custom commands -- common operators
% --------------------------------------------------------------------------
\DeclareMathOperator{\dom}{dom}
\DeclareMathOperator{\epi}{epi}
\DeclareMathOperator{\conv}{conv}
\DeclareMathOperator{\relint}{relint}
\DeclareMathOperator{\aff}{aff}
\DeclareMathOperator{\cl}{cl}
\DeclareMathOperator{\interior}{int}
\DeclareMathOperator{\tr}{tr}
\DeclareMathOperator{\diag}{diag}
\DeclareMathOperator{\rank}{rank}
\DeclareMathOperator{\sgn}{sgn}
\DeclareMathOperator{\supp}{supp}
\DeclareMathOperator{\argmin}{arg\,min}
\DeclareMathOperator{\argmax}{arg\,max}
\DeclareMathOperator{\prox}{prox}
\DeclareMathOperator{\dist}{dist}
\DeclareMathOperator{\diam}{diam}
\DeclareMathOperator{\Span}{span}

% --------------------------------------------------------------------------
% Custom commands -- linear algebra operators
% --------------------------------------------------------------------------
\DeclareMathOperator{\nullity}{nullity}
\DeclareMathOperator{\range}{range}
\DeclareMathOperator{\nullspace}{null}
\DeclareMathOperator{\proj}{proj}

% --------------------------------------------------------------------------
% Custom commands -- measure theory
% --------------------------------------------------------------------------
\newcommand{\powerset}{\mathcal{P}}
\newcommand{\Borel}{\mathcal{B}}

\pagestyle{fancy}
\fancyhf{}
\fancyhead[L]{\textsc{Lesson 9: Probability Spaces and Conditioning}}
\fancyhead[R]{\thepage}
\title{Lesson 9: Probability Spaces and Conditioning}
\author{AI/ML Mathematics}
\date{}
\begin{document}
\maketitle
\tableofcontents
\newpage

%% ============================================================
%% SECTION 1: MOTIVATION AND OVERVIEW
%% ============================================================

\section{Motivation and Overview}
\label{sec:motivation}

Consider a reinforcement learning agent navigating a stochastic gridworld.
At each time step the agent selects an action, but the resulting state
transition is random: a ``move north'' command might succeed with probability
$0.8$ and slip sideways with probability $0.1$ in each perpendicular
direction. To compute expected returns, compare policies, or update beliefs
after observing a trajectory, the agent needs a rigorous language for
quantifying uncertainty. What does it mean for an event to have probability
$0.8$? How should the agent revise its beliefs about a hidden parameter after
collecting data? If two noise sources affect the system, when can their
contributions be analyzed separately?

The tools developed in this lesson will let us formalize these questions.
We construct probability spaces via the Kolmogorov axioms, derive the
fundamental rules for computing with probabilities (complement, monotonicity,
inclusion-exclusion, continuity), and develop conditional probability, Bayes'
theorem, and the concept of independence. These results form the
prerequisite language for every probabilistic statement in reinforcement
learning and machine learning.

\begin{rlconnection}[title={Why Probability Spaces Matter for RL/ML}]
\begin{itemize}[nosep]
    \item \textbf{Markov decision processes}: the transition kernel
    $p(s' \mid s, a)$ is a conditional probability; the Markov property is a
    conditional independence statement.
    \item \textbf{Bayesian inference}: posterior updates in Thompson sampling
    and Bayesian neural networks apply Bayes' theorem at every step.
    \item \textbf{Monte Carlo methods}: REINFORCE and Monte Carlo tree search
    estimate expectations by sampling from probability spaces.
    \item \textbf{PAC learning bounds}: generalization guarantees in
    supervised learning rely on the union bound, inclusion-exclusion, and
    independence of training samples.
\end{itemize}
\end{rlconnection}

\paragraph{Prerequisites.}
Basic set theory: unions, intersections, complements, De Morgan's laws
(Lesson~1). Familiarity with sequences and limits of real numbers
(Lessons~2--3). No linear algebra is required for this lesson.

\paragraph{Learning objectives.}
After completing this lesson, the student will be able to:
\begin{itemize}
    \item \emph{Define} sample spaces, events, and probability measures via
    the Kolmogorov axioms.
    \item \emph{Derive} basic properties of probability (complement rule,
    monotonicity, inclusion-exclusion) from the axioms.
    \item \emph{Prove} the continuity of probability for monotone sequences
    of events.
    \item \emph{Compute} conditional probabilities and apply the
    multiplication rule, law of total probability, and Bayes' theorem.
    \item \emph{Distinguish} between pairwise and mutual independence and
    \emph{construct} examples separating the two notions.
    \item \emph{Apply} Bayes' theorem to belief-update problems arising in
    machine learning.
\end{itemize}

%% ============================================================
%% SECTION 2: REFERENCES
%% ============================================================

\section{References}
\label{sec:references}

\begin{itemize}
    \item \citet[Chapters 1--2]{ross2019} -- Combinatorial analysis, axioms
    of probability, conditional probability and independence.
    \item \citet[Chapters 1--2]{blitzstein2019} -- Probability and counting,
    conditional probability with extensive examples.
    \item \citet[Chapter 1]{feller1968} -- Classical introduction to sample
    spaces and probability axioms.
    \item \citet[Chapter 1]{grimmett2020} -- Events and their probabilities,
    with a path toward measure-theoretic foundations.
    \item \citet[Chapters 1--2]{bertsekas2008} -- Sample space, probability
    axioms, conditioning, and independence with engineering applications.
\end{itemize}

%% ============================================================
%% SECTION 3: PROBABILITY SPACES
%% ============================================================

\section{Probability Spaces}
\label{sec:prob_spaces}

We begin by defining the mathematical objects that model random experiments.

\begin{definition}[Sample Space, Outcome, Event]
\label{def:sample_space}
A \emph{sample space} $\Omega$ is a non-empty set. Each element
$\omega \in \Omega$ is called an \emph{outcome}. An \emph{event} is a
subset $A \subseteq \Omega$.
\end{definition}

\begin{remark}
\label{rem:power_set}
In this lesson we take the collection of all events to be the power set
$2^\Omega$ (the set of all subsets of $\Omega$). For uncountable sample
spaces, one must restrict to a $\sigma$-algebra
$\mathcal{F} \subseteq 2^\Omega$ to avoid measure-theoretic paradoxes;
this refinement is developed in Phase~01 (Lesson~1, $\sigma$-algebras).
\end{remark}

\begin{definition}[Set Operations on Events]
\label{def:set_ops}
Let $A, B \subseteq \Omega$.
\begin{enumerate}[nosep]
    \item \emph{Union:}
    $A \cup B = \{\omega \in \Omega : \omega \in A \text{ or }
    \omega \in B\}$ (``$A$ or $B$ occurs'').
    \item \emph{Intersection:}
    $A \cap B = \{\omega \in \Omega : \omega \in A \text{ and }
    \omega \in B\}$ (``both $A$ and $B$ occur'').
    \item \emph{Complement:}
    $A^c = \Omega \setminus A = \{\omega \in \Omega : \omega \notin A\}$
    (``$A$ does not occur'').
    \item \emph{Disjoint:} $A$ and $B$ are \emph{disjoint} (or
    \emph{mutually exclusive}) if $A \cap B = \emptyset$.
\end{enumerate}
\end{definition}

\begin{example}[Die Roll]
\label{ex:die_roll}
Roll a fair six-sided die. The sample space is
$\Omega = \{1, 2, 3, 4, 5, 6\}$. Let $A = \{2, 4, 6\}$ (``even'') and
$B = \{1, 2, 3\}$ (``at most 3''). Then $A \cup B = \{1, 2, 3, 4, 6\}$,
$A \cap B = \{2\}$, and $A^c = \{1, 3, 5\}$.
\end{example}

\subsection{The Kolmogorov Axioms}

\begin{definition}[Probability Measure -- Kolmogorov Axioms]
\label{def:kolmogorov}
A \emph{probability measure} (or simply a \emph{probability}) on a sample
space $\Omega$ is a function $\Prob \colon 2^\Omega \to \R$ satisfying:
\begin{enumerate}[nosep,label=(K\arabic*)]
    \item \textbf{Non-negativity:} $\Prob(A) \geq 0$ for every event $A$.
    \item \textbf{Normalization:} $\Prob(\Omega) = 1$.
    \item \textbf{Countable additivity:} If $A_1, A_2, \ldots$ are pairwise
    disjoint events, then
    \[
        \Prob\!\left(\bigcup_{n=1}^\infty A_n\right)
        = \sum_{n=1}^\infty \Prob(A_n).
    \]
\end{enumerate}
The pair $(\Omega, \Prob)$ is called a \emph{probability space}.
\end{definition}

\begin{remark}
\label{rem:finite_additivity}
Countable additivity (K3) implies \emph{finite additivity}: if
$A_1, \ldots, A_n$ are pairwise disjoint, set $A_k = \emptyset$ for
$k > n$. Once we prove $\Prob(\emptyset) = 0$ below, the tail terms
vanish and $\Prob(\bigcup_{i=1}^n A_i) = \sum_{i=1}^n \Prob(A_i)$.
\end{remark}

\begin{example}[Uniform Probability on a Finite Set]
\label{ex:uniform}
Let $\Omega$ be a finite set with $|\Omega| = N$. Define
$\Prob(A) = |A|/N$ for every $A \subseteq \Omega$. Then (K1) holds
because $|A| \geq 0$; (K2) holds because $|\Omega|/N = 1$; (K3) holds
because for disjoint sets $|A \cup B| = |A| + |B|$ (and only finitely many
events are non-empty, so the infinite sum reduces to a finite one).
\end{example}

\subsection{Basic Properties}

We now derive the fundamental consequences of the Kolmogorov axioms.

\begin{proposition}[Complement Rule]
\label{prop:complement}
For any event $A$, $\Prob(A^c) = 1 - \Prob(A)$.
\end{proposition}

\begin{proof}
Since $A$ and $A^c$ are disjoint and $A \cup A^c = \Omega$, countable
additivity (K3) (applied to two events, with the remaining sets empty)
gives $\Prob(\Omega) = \Prob(A) + \Prob(A^c)$. By normalization (K2),
$\Prob(\Omega) = 1$, so $\Prob(A^c) = 1 - \Prob(A)$.
\end{proof}

\begin{corollary}[$\Prob(\emptyset) = 0$]
\label{cor:empty_set}
$\Prob(\emptyset) = 0$.
\end{corollary}

\begin{proof}
Since $\emptyset = \Omega^c$, the complement rule gives
$\Prob(\emptyset) = 1 - \Prob(\Omega) = 1 - 1 = 0$.
\end{proof}

\begin{proposition}[Monotonicity]
\label{prop:monotonicity}
If $A \subseteq B$, then $\Prob(A) \leq \Prob(B)$.
\end{proposition}

\begin{proof}
Write $B = A \cup (B \setminus A)$, where $A$ and $B \setminus A$ are
disjoint. By finite additivity,
$\Prob(B) = \Prob(A) + \Prob(B \setminus A)$. Since
$\Prob(B \setminus A) \geq 0$ by (K1), we conclude
$\Prob(B) \geq \Prob(A)$.
\end{proof}

\begin{example}[Applying Monotonicity]
\label{ex:monotonicity_app}
Roll a fair die. Let $A = \{6\}$ and $B = \{4, 5, 6\}$. Since
$A \subseteq B$, monotonicity guarantees $\Prob(A) \leq \Prob(B)$.
Indeed, $1/6 \leq 3/6 = 1/2$.
\end{example}

\begin{keyresult}[title={Inclusion-Exclusion Principle}]
The inclusion-exclusion formula generalizes finite additivity to
overlapping events. For two events it corrects for double-counting;
for $n$ events it alternates between over- and under-counting.
\end{keyresult}

\begin{theorem}[Inclusion-Exclusion for Two Events]
\label{thm:ie_two}
For any events $A$ and $B$,
\[
    \Prob(A \cup B) = \Prob(A) + \Prob(B) - \Prob(A \cap B).
\]
\end{theorem}

\begin{proof}
Write $A \cup B = A \cup (B \setminus A)$, where $A$ and $B \setminus A$
are disjoint. By finite additivity,
\begin{equation}
\label{eq:ie_step1}
    \Prob(A \cup B) = \Prob(A) + \Prob(B \setminus A).
\end{equation}
Similarly, $B = (A \cap B) \cup (B \setminus A)$ with
$A \cap B$ and $B \setminus A$ disjoint, so
\begin{equation}
\label{eq:ie_step2}
    \Prob(B) = \Prob(A \cap B) + \Prob(B \setminus A),
\end{equation}
which gives $\Prob(B \setminus A) = \Prob(B) - \Prob(A \cap B)$.
Substituting into \eqref{eq:ie_step1}:
\[
    \Prob(A \cup B) = \Prob(A) + \Prob(B) - \Prob(A \cap B). \qedhere
\]
\end{proof}

\begin{example}[Inclusion-Exclusion Applied]
\label{ex:ie_applied}
Draw one card from a standard 52-card deck. Let $A$ = ``the card is a
heart'' and $B$ = ``the card is a face card (J, Q, K).'' Then
$\Prob(A) = 13/52$, $\Prob(B) = 12/52$, and
$\Prob(A \cap B) = 3/52$ (the three heart face cards). By
inclusion-exclusion,
$\Prob(A \cup B) = 13/52 + 12/52 - 3/52 = 22/52 = 11/26$.
\end{example}

\begin{theorem}[Inclusion-Exclusion for $n$ Events]
\label{thm:ie_n}
For events $A_1, \ldots, A_n$,
\[
    \Prob\!\left(\bigcup_{i=1}^n A_i\right)
    = \sum_{k=1}^n (-1)^{k+1}
    \sum_{1 \leq i_1 < \cdots < i_k \leq n}
    \Prob(A_{i_1} \cap \cdots \cap A_{i_k}).
\]
\end{theorem}

\begin{proof}
We proceed by induction on $n$.

\emph{Base case} ($n = 1$): $\Prob(A_1) = \Prob(A_1)$, which is trivially
the formula with a single term.

\emph{Base case} ($n = 2$): This is Theorem~\ref{thm:ie_two}.

\emph{Inductive step.} Assume the formula holds for $n - 1$ events. Write
\[
    \bigcup_{i=1}^n A_i
    = \left(\bigcup_{i=1}^{n-1} A_i\right) \cup A_n.
\]
By Theorem~\ref{thm:ie_two},
\begin{equation}
\label{eq:ie_induction}
    \Prob\!\left(\bigcup_{i=1}^n A_i\right)
    = \Prob\!\left(\bigcup_{i=1}^{n-1} A_i\right) + \Prob(A_n)
    - \Prob\!\left(\left(\bigcup_{i=1}^{n-1} A_i\right) \cap A_n\right).
\end{equation}
By distributivity,
$\left(\bigcup_{i=1}^{n-1} A_i\right) \cap A_n
= \bigcup_{i=1}^{n-1} (A_i \cap A_n)$.
Apply the inductive hypothesis to the first term in \eqref{eq:ie_induction}
(a union of $n-1$ events $A_1, \ldots, A_{n-1}$) and to the third term
(a union of $n-1$ events $A_1 \cap A_n, \ldots, A_{n-1} \cap A_n$).

The first term contributes the alternating sum over all subsets of
$\{1, \ldots, n-1\}$. Adding $\Prob(A_n)$ accounts for the singleton
$\{n\}$. Subtracting the inductive hypothesis applied to
$\bigcup_{i=1}^{n-1}(A_i \cap A_n)$ contributes exactly the terms
involving $A_n$ in subsets of size $\geq 2$, each with sign
$(-1)^{k+1}$ (the extra intersection with $A_n$ increases the subset
size by one and flips the sign). Collecting all terms yields the formula
for $n$ events.
\end{proof}

\begin{warning}[title={Union Bound as a Corollary}]
A common shortcut is the \emph{union bound} (Boole's inequality):
$\Prob(\bigcup_{i=1}^n A_i) \leq \sum_{i=1}^n \Prob(A_i)$.
This follows from inclusion-exclusion by dropping the
alternating correction terms (or directly from monotonicity and finite
additivity). In machine learning, the union bound is the workhorse
behind PAC learning arguments and uniform convergence bounds. Be aware
that it can be very loose when events overlap significantly.
\end{warning}

\subsection{Continuity of Probability}

The following theorem connects the axioms to limits of sequences of events
and is essential for analyzing asymptotic properties.

\begin{theorem}[Continuity of Probability]
\label{thm:continuity}
\begin{enumerate}[nosep,label=(\alph*)]
    \item \textbf{Continuity from below.} If
    $A_1 \subseteq A_2 \subseteq A_3 \subseteq \cdots$ (increasing), then
    \[
        \Prob\!\left(\bigcup_{n=1}^\infty A_n\right)
        = \lim_{n \to \infty} \Prob(A_n).
    \]
    \item \textbf{Continuity from above.} If
    $A_1 \supseteq A_2 \supseteq A_3 \supseteq \cdots$ (decreasing), then
    \[
        \Prob\!\left(\bigcap_{n=1}^\infty A_n\right)
        = \lim_{n \to \infty} \Prob(A_n).
    \]
\end{enumerate}
\end{theorem}

\begin{proof}
\textbf{(a)} Define $B_1 = A_1$ and $B_n = A_n \setminus A_{n-1}$ for
$n \geq 2$. The events $B_1, B_2, \ldots$ are pairwise disjoint: if
$i < j$, then $B_i \subseteq A_i \subseteq A_{j-1}$, but
$B_j = A_j \setminus A_{j-1}$ is disjoint from $A_{j-1}$, so
$B_i \cap B_j = \emptyset$. Moreover,
$\bigcup_{n=1}^\infty B_n = \bigcup_{n=1}^\infty A_n$. By countable
additivity (K3),
\[
    \Prob\!\left(\bigcup_{n=1}^\infty A_n\right)
    = \sum_{n=1}^\infty \Prob(B_n).
\]
The partial sums satisfy
$\sum_{k=1}^n \Prob(B_k)
= \Prob\!\left(\bigcup_{k=1}^n B_k\right)
= \Prob(A_n)$,
where the first equality uses finite additivity (the $B_k$ are disjoint)
and the second uses
$\bigcup_{k=1}^n B_k = A_n$ (proved by induction: $B_1 = A_1$, and
adding $B_n = A_n \setminus A_{n-1}$ to $A_{n-1}$ gives $A_n$).
Therefore
\[
    \Prob\!\left(\bigcup_{n=1}^\infty A_n\right)
    = \lim_{n \to \infty} \sum_{k=1}^n \Prob(B_k)
    = \lim_{n \to \infty} \Prob(A_n).
\]

\textbf{(b)} If $A_1 \supseteq A_2 \supseteq \cdots$, then the complements
$A_1^c \subseteq A_2^c \subseteq \cdots$ form an increasing sequence. By
part (a),
\[
    \Prob\!\left(\bigcup_{n=1}^\infty A_n^c\right)
    = \lim_{n \to \infty} \Prob(A_n^c)
    = \lim_{n \to \infty} (1 - \Prob(A_n)).
\]
By De Morgan's law,
$\bigcup_{n=1}^\infty A_n^c = \left(\bigcap_{n=1}^\infty A_n\right)^c$,
so the left side equals
$1 - \Prob\!\left(\bigcap_{n=1}^\infty A_n\right)$. Equating and
simplifying:
$\Prob\!\left(\bigcap_{n=1}^\infty A_n\right)
= \lim_{n \to \infty} \Prob(A_n)$.
\end{proof}

\begin{example}[Continuity from Below]
\label{ex:continuity_below}
Let $\Omega = (0, 1]$ with the uniform distribution (Lebesgue measure).
Define $A_n = (0, 1 - 1/n]$ for $n \geq 2$. Then $A_2 \subseteq A_3
\subseteq \cdots$ and $\bigcup_{n=2}^\infty A_n = (0, 1)$. By continuity
from below,
$\Prob((0, 1)) = \lim_{n \to \infty} \Prob(A_n) = \lim_{n \to \infty}
(1 - 1/n) = 1$.
This confirms that $\Prob(\{1\}) = 0$ for a continuous uniform
distribution.
\end{example}

\begin{rlconnection}[title={Continuity and Convergence in RL}]
Continuity from above is essential for ``almost sure'' convergence
statements in RL. The event ``the value estimate $V_t$ converges to the
true value $V^\pi$'' is expressed as
$\bigcap_{k=1}^\infty \bigcup_{N=1}^\infty
\bigcap_{t=N}^\infty \{|V_t - V^\pi| < 1/k\}$. Computing the
probability of such nested limits requires the continuity of $\Prob$
along monotone sequences.
\end{rlconnection}

We have established the axiomatic framework and its basic consequences.
The next section introduces conditional probability, which allows us to
update probabilities given partial information.

%% ============================================================
%% SECTION 4: CONDITIONAL PROBABILITY
%% ============================================================

\section{Conditional Probability}
\label{sec:cond_prob}

In many situations we have partial information about the outcome of an
experiment, and we wish to update our probability assignments accordingly.
Conditional probability formalizes this update.

\begin{definition}[Conditional Probability]
\label{def:cond_prob}
Let $(\Omega, \Prob)$ be a probability space and let $B$ be an event with
$\Prob(B) > 0$. The \emph{conditional probability of $A$ given $B$} is
\[
    \Prob(A \mid B) = \frac{\Prob(A \cap B)}{\Prob(B)}.
\]
\end{definition}

\begin{remark}
\label{rem:cond_is_prob}
For a fixed event $B$ with $\Prob(B) > 0$, the function
$A \mapsto \Prob(A \mid B)$ is itself a probability measure on $\Omega$.
Verification: (K1) $\Prob(A \mid B) = \Prob(A \cap B)/\Prob(B) \geq 0$;
(K2) $\Prob(\Omega \mid B) = \Prob(B)/\Prob(B) = 1$; (K3) if
$A_1, A_2, \ldots$ are pairwise disjoint, then
$\Prob(\bigcup_n A_n \mid B)
= \Prob((\bigcup_n A_n) \cap B)/\Prob(B)
= \Prob(\bigcup_n (A_n \cap B))/\Prob(B)
= \sum_n \Prob(A_n \cap B)/\Prob(B)
= \sum_n \Prob(A_n \mid B)$.
\end{remark}

\begin{example}[Conditional Probability with Dice]
\label{ex:cond_dice}
Roll two fair dice. Let $A$ = ``the sum is 8'' and $B$ = ``the first die
shows 3.'' The sample space has $|\Omega| = 36$ equally likely outcomes.
We have $A \cap B = \{(3, 5)\}$, so $\Prob(A \cap B) = 1/36$, and
$\Prob(B) = 6/36 = 1/6$. Thus
$\Prob(A \mid B) = (1/36)/(1/6) = 1/6$.
\end{example}

\subsection{Multiplication Rule}

The multiplication rule extends the definition of conditional probability
to chains of events.

\begin{theorem}[Multiplication Rule]
\label{thm:mult_rule}
For events $A_1, \ldots, A_n$ with
$\Prob(A_1 \cap \cdots \cap A_{n-1}) > 0$,
\[
    \Prob(A_1 \cap \cdots \cap A_n)
    = \Prob(A_1) \, \Prob(A_2 \mid A_1) \,
    \Prob(A_3 \mid A_1 \cap A_2) \cdots
    \Prob(A_n \mid A_1 \cap \cdots \cap A_{n-1}).
\]
\end{theorem}

\begin{proof}
We proceed by induction on $n$.

\emph{Base case} ($n = 2$): By definition of conditional probability,
$\Prob(A_1 \cap A_2)
= \Prob(A_1) \, \Prob(A_2 \mid A_1)$,
since $\Prob(A_2 \mid A_1) = \Prob(A_1 \cap A_2)/\Prob(A_1)$ and
$\Prob(A_1) \geq \Prob(A_1 \cap A_2) > 0$ by monotonicity (as
$A_1 \cap A_2 \subseteq A_1 \cap \cdots \cap A_{n-1}$ and
$\Prob(A_1 \cap \cdots \cap A_{n-1}) > 0$ implies
$\Prob(A_1) > 0$).

\emph{Inductive step.} Assume the formula holds for $n - 1$ events.
Since $\Prob(A_1 \cap \cdots \cap A_{n-1}) > 0$, we can write
\begin{align*}
    \Prob(A_1 \cap \cdots \cap A_n)
    &= \Prob(A_1 \cap \cdots \cap A_{n-1}) \cdot
       \Prob(A_n \mid A_1 \cap \cdots \cap A_{n-1}) \\
    &= \Bigl[\Prob(A_1) \, \Prob(A_2 \mid A_1) \cdots
       \Prob(A_{n-1} \mid A_1 \cap \cdots \cap A_{n-2})\Bigr] \\
    &\qquad \cdot \Prob(A_n \mid A_1 \cap \cdots \cap A_{n-1}),
\end{align*}
where the first line uses the definition of conditional probability and
the second applies the inductive hypothesis to the $n - 1$ events
$A_1, \ldots, A_{n-1}$.
\end{proof}

\begin{example}[Drawing Without Replacement]
\label{ex:draw_without_replacement}
An urn contains 5 red and 3 blue balls. We draw 3 balls without
replacement. Let $A_i$ = ``the $i$-th ball drawn is red.'' By the
multiplication rule,
\begin{align*}
    \Prob(A_1 \cap A_2 \cap A_3)
    &= \Prob(A_1) \, \Prob(A_2 \mid A_1) \,
       \Prob(A_3 \mid A_1 \cap A_2) \\
    &= \frac{5}{8} \cdot \frac{4}{7} \cdot \frac{3}{6}
    = \frac{60}{336} = \frac{5}{28}.
\end{align*}
\end{example}

\subsection{Law of Total Probability}

\begin{definition}[Partition]
\label{def:partition}
Events $B_1, B_2, \ldots, B_n$ form a \emph{partition} of $\Omega$ if
they are pairwise disjoint and $\bigcup_{i=1}^n B_i = \Omega$.
\end{definition}

\begin{theorem}[Law of Total Probability]
\label{thm:total_prob}
Let $B_1, B_2, \ldots, B_n$ be a partition of $\Omega$ with
$\Prob(B_i) > 0$ for all $i$. Then for any event $A$,
\[
    \Prob(A) = \sum_{i=1}^n \Prob(A \mid B_i) \, \Prob(B_i).
\]
\end{theorem}

\begin{proof}
Since the $B_i$ partition $\Omega$,
\[
    A = A \cap \Omega = A \cap \left(\bigcup_{i=1}^n B_i\right)
    = \bigcup_{i=1}^n (A \cap B_i).
\]
The events $A \cap B_1, \ldots, A \cap B_n$ are pairwise disjoint (since
the $B_i$ are). By finite additivity,
\[
    \Prob(A) = \sum_{i=1}^n \Prob(A \cap B_i)
    = \sum_{i=1}^n \Prob(A \mid B_i) \, \Prob(B_i),
\]
where the last equality uses
$\Prob(A \cap B_i) = \Prob(A \mid B_i) \, \Prob(B_i)$ from the definition
of conditional probability.
\end{proof}

\begin{example}[Medical Test -- Total Probability]
\label{ex:medical_total}
A disease affects 1\% of the population. A diagnostic test has a
sensitivity (true positive rate) of 95\% and a specificity (true negative
rate) of 90\%. Let $D$ = ``has disease'' and $T$ = ``tests positive.''
The partition is $\{D, D^c\}$ with $\Prob(D) = 0.01$ and
$\Prob(D^c) = 0.99$. By the law of total probability,
\begin{align*}
    \Prob(T) &= \Prob(T \mid D) \, \Prob(D)
               + \Prob(T \mid D^c) \, \Prob(D^c) \\
             &= 0.95 \times 0.01 + 0.10 \times 0.99 \\
             &= 0.0095 + 0.099 = 0.1085.
\end{align*}
Despite the test's high sensitivity, most positive results are false
positives because the disease is rare.
\end{example}

\subsection{Bayes' Theorem}

\begin{keyresult}[title={Bayes' Theorem -- The Belief Update Rule}]
Bayes' theorem inverts the direction of conditioning: given a partition
into ``hypotheses'' $B_1, \ldots, B_n$ with prior probabilities
$\Prob(B_j)$, observing event $A$ updates each hypothesis to the
posterior $\Prob(B_j \mid A)$.
\end{keyresult}

\begin{theorem}[Bayes' Theorem]
\label{thm:bayes}
Let $B_1, \ldots, B_n$ be a partition of $\Omega$ with
$\Prob(B_i) > 0$ for all $i$, and let $A$ be an event with
$\Prob(A) > 0$. Then for each $j = 1, \ldots, n$,
\[
    \Prob(B_j \mid A)
    = \frac{\Prob(A \mid B_j) \, \Prob(B_j)}
           {\sum_{i=1}^n \Prob(A \mid B_i) \, \Prob(B_i)}.
\]
\end{theorem}

\begin{proof}
By the definition of conditional probability,
\[
    \Prob(B_j \mid A) = \frac{\Prob(A \cap B_j)}{\Prob(A)}.
\]
The numerator equals
$\Prob(A \mid B_j) \, \Prob(B_j)$
by the definition of $\Prob(A \mid B_j)$. The denominator equals
$\sum_{i=1}^n \Prob(A \mid B_i) \, \Prob(B_i)$
by the law of total probability (Theorem~\ref{thm:total_prob}).
Substituting gives the result.
\end{proof}

\begin{example}[Medical Test -- Bayes' Theorem]
\label{ex:medical_bayes}
Continuing Example~\ref{ex:medical_total}, suppose a patient tests
positive. What is the probability they actually have the disease?
By Bayes' theorem,
\[
    \Prob(D \mid T)
    = \frac{\Prob(T \mid D) \, \Prob(D)}{\Prob(T)}
    = \frac{0.95 \times 0.01}{0.1085}
    \approx 0.0876.
\]
Despite a positive test result, the probability of disease is only about
8.8\%. This illustrates the \emph{base rate fallacy}: ignoring the low
prior $\Prob(D)$ leads to grossly overestimating $\Prob(D \mid T)$.
\end{example}

\begin{example}[Urn Problem -- Bayesian Update]
\label{ex:urn_bayes}
There are two urns. Urn 1 has 7 red and 3 blue balls; Urn 2 has 4 red
and 6 blue balls. One urn is chosen uniformly at random, and a ball is
drawn. Let $U_i$ = ``Urn $i$ was chosen'' and $R$ = ``a red ball is
drawn.'' The prior is $\Prob(U_1) = \Prob(U_2) = 1/2$. We have
$\Prob(R \mid U_1) = 7/10$ and $\Prob(R \mid U_2) = 4/10$.

Given that a red ball was drawn, Bayes' theorem gives
\[
    \Prob(U_1 \mid R)
    = \frac{(7/10)(1/2)}{(7/10)(1/2) + (4/10)(1/2)}
    = \frac{7/20}{11/20}
    = \frac{7}{11} \approx 0.636.
\]
The observation of red shifts our belief toward Urn~1.
\end{example}

\begin{rlconnection}[title={Bayes' Theorem in RL/ML}]
\begin{itemize}[nosep]
    \item In \textbf{Bayesian reinforcement learning}, the agent maintains a
    prior over transition models and updates it via Bayes' theorem after
    each observed transition.
    \item In \textbf{POMDPs}, the belief state
    $b'(s') \propto \Prob(o \mid s') \sum_s \Prob(s' \mid s, a) \, b(s)$
    is a direct application of the law of total probability and Bayes'
    theorem.
    \item In \textbf{naive Bayes classifiers}, the class posterior is
    computed via Bayes' theorem with a conditional independence assumption
    on features.
\end{itemize}
\end{rlconnection}

We now turn to the concept of independence, which characterizes when
knowledge of one event provides no information about another.

%% ============================================================
%% SECTION 5: INDEPENDENCE
%% ============================================================

\section{Independence}
\label{sec:independence}

\begin{definition}[Independence of Two Events]
\label{def:indep_two}
Events $A$ and $B$ are \emph{independent} if
\[
    \Prob(A \cap B) = \Prob(A) \, \Prob(B).
\]
\end{definition}

\begin{remark}
\label{rem:indep_conditional}
When $\Prob(B) > 0$, independence is equivalent to
$\Prob(A \mid B) = \Prob(A)$: learning that $B$ occurred does not change
the probability of $A$. This follows directly from the definition:
$\Prob(A \mid B) = \Prob(A \cap B)/\Prob(B) = \Prob(A)\Prob(B)/\Prob(B)
= \Prob(A)$.
\end{remark}

\begin{example}[Independent Events]
\label{ex:indep_coins}
Toss a fair coin twice. Let $A$ = ``first toss is heads'' and
$B$ = ``second toss is heads.'' The sample space is
$\Omega = \{HH, HT, TH, TT\}$ with each outcome having probability
$1/4$. Then $\Prob(A) = 1/2$, $\Prob(B) = 1/2$, and
$\Prob(A \cap B) = \Prob(\{HH\}) = 1/4 = (1/2)(1/2)
= \Prob(A)\Prob(B)$. So $A$ and $B$ are independent.
\end{example}

\begin{warning}[title={Independent $\neq$ Disjoint}]
A common mistake is to confuse independence with disjointness. If $A$ and
$B$ are disjoint and both have positive probability, then
$\Prob(A \cap B) = 0 \neq \Prob(A)\Prob(B) > 0$, so they are
\emph{not} independent. Disjoint events are maximally dependent:
knowing one occurred tells you the other did not.
\end{warning}

\begin{proposition}[Independence of Complements]
\label{prop:indep_complement}
If $A$ and $B$ are independent, then so are $A$ and $B^c$, $A^c$ and $B$,
and $A^c$ and $B^c$.
\end{proposition}

\begin{proof}
We prove that $A$ and $B^c$ are independent; the other cases follow by
symmetry and iteration. Since $A = (A \cap B) \cup (A \cap B^c)$ with
the two sets disjoint,
\[
    \Prob(A \cap B^c) = \Prob(A) - \Prob(A \cap B)
    = \Prob(A) - \Prob(A)\Prob(B)
    = \Prob(A)(1 - \Prob(B))
    = \Prob(A)\Prob(B^c).
\]
For $A^c$ and $B^c$: since $A$ and $B^c$ are independent (just proved),
applying the same argument with $A$ replaced by $A^c$ gives
$\Prob(A^c \cap B^c) = \Prob(A^c)\Prob(B^c)$.
\end{proof}

\subsection{Mutual Independence}

\begin{definition}[Mutual Independence]
\label{def:mutual_indep}
Events $A_1, \ldots, A_n$ are \emph{mutually independent} if for every
subset $\{i_1, \ldots, i_k\} \subseteq \{1, \ldots, n\}$ with
$k \geq 2$,
\[
    \Prob(A_{i_1} \cap \cdots \cap A_{i_k})
    = \Prob(A_{i_1}) \cdots \Prob(A_{i_k}).
\]
This requires the product formula to hold for \emph{all}
$\binom{n}{2} + \binom{n}{3} + \cdots + \binom{n}{n} = 2^n - n - 1$
subsets of size $\geq 2$, not merely for the $\binom{n}{2}$ pairs.
\end{definition}

\begin{remark}
\label{rem:pairwise_def}
Events $A_1, \ldots, A_n$ are \emph{pairwise independent} if every pair
is independent:
$\Prob(A_i \cap A_j) = \Prob(A_i)\Prob(A_j)$
for all $i \neq j$. Mutual independence implies pairwise independence
(take $k = 2$ in the definition), but the converse is false.
\end{remark}

\begin{warning}[title={Pairwise $\neq$ Mutual Independence}]
Pairwise independence does not imply mutual independence. The product
rule can hold for every pair yet fail for a triple (or larger subset).
This distinction matters in practice: many concentration inequalities
(e.g., the variance of a sum) require only pairwise independence, while
others (e.g., the strong law of large numbers) require mutual
independence.
\end{warning}

\begin{example}[Pairwise but Not Mutually Independent]
\label{ex:pairwise_not_mutual}
Toss two fair coins. The sample space is
$\Omega = \{HH, HT, TH, TT\}$, each outcome with probability $1/4$.
Define:
\begin{itemize}[nosep]
    \item $A$ = ``the first coin is heads'' $= \{HH, HT\}$,
    \item $B$ = ``the second coin is heads'' $= \{HH, TH\}$,
    \item $C$ = ``exactly one head'' $= \{HT, TH\}$.
\end{itemize}
Each event has probability $1/2$. We check all pairs:
\begin{itemize}[nosep]
    \item $A \cap B = \{HH\}$, so
    $\Prob(A \cap B) = 1/4 = (1/2)(1/2) = \Prob(A)\Prob(B)$. \checkmark
    \item $A \cap C = \{HT\}$, so
    $\Prob(A \cap C) = 1/4 = (1/2)(1/2) = \Prob(A)\Prob(C)$. \checkmark
    \item $B \cap C = \{TH\}$, so
    $\Prob(B \cap C) = 1/4 = (1/2)(1/2) = \Prob(B)\Prob(C)$. \checkmark
\end{itemize}
All three pairs are independent. However,
$A \cap B \cap C = \emptyset$ (both heads is not exactly one head), so
\[
    \Prob(A \cap B \cap C) = 0
    \neq \frac{1}{8}
    = \Prob(A)\Prob(B)\Prob(C).
\]
Therefore $A$, $B$, $C$ are pairwise independent but \emph{not}
mutually independent. Knowing $A$ and $B$ both occurred completely
determines $C$ (it cannot occur), even though each pair individually
reveals nothing.
\end{example}

\begin{example}[Mutual Independence in Repeated Trials]
\label{ex:mutual_indep_trials}
Toss a fair coin three times independently. The sample space has
$|\Omega| = 8$ equally likely outcomes. Let $A_i$ = ``the $i$-th toss
is heads'' for $i = 1, 2, 3$. Then $\Prob(A_i) = 1/2$ for each $i$.
For pairs:
$\Prob(A_i \cap A_j) = 2/8 = 1/4 = (1/2)(1/2)$ for each pair
$i \neq j$ (two of the eight outcomes have heads on both specified
positions). For the triple:
$\Prob(A_1 \cap A_2 \cap A_3) = 1/8 = (1/2)^3$.
All $2^3 - 3 - 1 = 4$ conditions are satisfied, confirming mutual
independence.
\end{example}

\begin{rlconnection}[title={Independence and the Markov Property}]
The Markov property is fundamentally a conditional independence statement:
the future evolution of the process is conditionally independent of the
past given the present state. Formally, in a Markov decision process,
$\Prob(S_{t+1} \mid S_0, A_0, \ldots, S_t, A_t) = \Prob(S_{t+1} \mid S_t, A_t)$.
Independence of noise terms across time steps justifies treating
transition randomness as i.i.d., which is central to the convergence
analysis of TD-learning and Q-learning.
\end{rlconnection}

%% ============================================================
%% SECTION 6: EXERCISES
%% ============================================================

\section{Exercises}
\label{sec:exercises}

\begin{exercise}[Bonferroni Inequalities]
\textnormal{(\(\ast\))} \\
Using inclusion-exclusion (Theorem~\ref{thm:ie_n}), prove the
\emph{first Bonferroni inequality}: for any events $A_1, \ldots, A_n$,
\[
    \Prob\!\left(\bigcup_{i=1}^n A_i\right)
    \leq \sum_{i=1}^n \Prob(A_i).
\]
Then prove the \emph{second Bonferroni inequality}:
\[
    \Prob\!\left(\bigcup_{i=1}^n A_i\right)
    \geq \sum_{i=1}^n \Prob(A_i)
    - \sum_{1 \leq i < j \leq n} \Prob(A_i \cap A_j).
\]
\textit{Hint:} For the first, use induction or monotonicity. For the
second, use inclusion-exclusion truncated at $k = 2$ and show the
remaining terms contribute non-negatively.
\end{exercise}

\begin{exercise}[Continuity from Above]
\textnormal{(\(\ast\))} \\
Let $A_1 \supseteq A_2 \supseteq \cdots$ with
$\Prob(A_n) = 1/(n+1)$ for all $n \geq 1$. Compute
$\Prob\!\left(\bigcap_{n=1}^\infty A_n\right)$ using the continuity of
probability.
\end{exercise}

\begin{exercise}[Sequential Bayes Updates]
\textnormal{(\(\ast\ast\))} \\
A factory has two machines. Machine~1 produces 60\% of items and has a
3\% defect rate; Machine~2 produces 40\% of items and has a 5\% defect
rate. An item is selected at random.
\begin{enumerate}[nosep,label=(\alph*)]
    \item Compute the probability that the item is defective.
    \item Given the item is defective, compute the probability it came
    from Machine~1.
    \item A second item is drawn (independently) and is also defective.
    Update the posterior for Machine~1 using Bayes' theorem applied
    to both observations.
\end{enumerate}
\end{exercise}

\begin{exercise}[Independence Structure in MDPs]
\textnormal{(\(\ast\ast\))} \\
Consider a Markov decision process with state space
$\mathcal{S} = \{s_1, s_2\}$, action space
$\mathcal{A} = \{a_1, a_2\}$, and deterministic policy
$\pi(s_1) = a_1$, $\pi(s_2) = a_2$. The transition probabilities are
$\Prob(S_{t+1} = s_1 \mid S_t = s_1, A_t = a_1) = 0.7$ and
$\Prob(S_{t+1} = s_1 \mid S_t = s_2, A_t = a_2) = 0.4$.
\begin{enumerate}[nosep,label=(\alph*)]
    \item Compute $\Prob(S_2 = s_1)$ given $\Prob(S_0 = s_1) = 1/2$
    using the law of total probability (applied twice).
    \item Are $\{S_0 = s_1\}$ and $\{S_2 = s_1\}$ independent? Compute
    $\Prob(S_0 = s_1, S_2 = s_1)$ and compare with
    $\Prob(S_0 = s_1) \Prob(S_2 = s_1)$.
\end{enumerate}
\end{exercise}

\begin{exercise}[Constructing Mutual Independence from Pairwise]
\textnormal{(\(\ast\ast\ast\))} \\
For $n = 4$, construct a finite probability space and events
$A_1, A_2, A_3, A_4$ that are pairwise independent but not mutually
independent. Verify all $\binom{4}{2} = 6$ pairwise conditions and
exhibit at least one subset of size $\geq 3$ for which the product
formula fails.

\textit{Hint:} Consider a sample space with $2^3 = 8$ equally likely
outcomes. Define each $A_i$ as a union of exactly 4 outcomes, chosen so
that every pair of events intersects in exactly 2 outcomes.
\end{exercise}

\begin{exercise}[Conditional Probability as a Probability Measure]
\textnormal{(\(\ast\ast\ast\))} \\
Let $(\Omega, \Prob)$ be a probability space and $B$ an event with
$\Prob(B) > 0$. Define $\Prob_B(A) = \Prob(A \mid B)$.
\begin{enumerate}[nosep,label=(\alph*)]
    \item Prove that $\Prob_B$ satisfies all three Kolmogorov axioms
    (this was stated in Remark~\ref{rem:cond_is_prob}; provide the full
    proof).
    \item Prove that if $A_1 \subseteq A_2 \subseteq \cdots$, then
    $\Prob_B(\bigcup_n A_n) = \lim_{n \to \infty} \Prob_B(A_n)$
    (continuity of conditional probability from below).
    \item Give an example showing that the continuity from above can fail
    if we allow $\Prob(B) = 0$, motivating the need for measure-theoretic
    conditional expectation (Phase~01).
\end{enumerate}
\end{exercise}

%% ============================================================
%% SECTION 7: SUMMARY
%% ============================================================

\section{Summary}
\label{sec:summary}

\begin{itemize}
    \item A \emph{probability space} $(\Omega, \Prob)$ assigns
    probabilities to events via three axioms: non-negativity,
    normalization ($\Prob(\Omega) = 1$), and countable additivity.
    \item The \emph{complement rule} ($\Prob(A^c) = 1 - \Prob(A)$) and
    \emph{monotonicity} ($A \subseteq B \Rightarrow \Prob(A) \leq
    \Prob(B)$) follow immediately from the axioms.
    \item \emph{Inclusion-exclusion} computes $\Prob(\bigcup A_i)$ via an
    alternating sum of intersection probabilities; truncating it gives the
    union bound $\Prob(\bigcup A_i) \leq \sum \Prob(A_i)$.
    \item \emph{Continuity of probability} extends the axioms to infinite
    limits: $\Prob(\bigcup A_n) = \lim \Prob(A_n)$ for increasing
    sequences and $\Prob(\bigcap A_n) = \lim \Prob(A_n)$ for decreasing
    sequences.
    \item \emph{Conditional probability}
    $\Prob(A \mid B) = \Prob(A \cap B)/\Prob(B)$ formalizes how partial
    information updates uncertainty; it is itself a probability measure.
    \item \emph{Bayes' theorem} inverts the direction of conditioning:
    $\Prob(B_j \mid A) \propto \Prob(A \mid B_j) \, \Prob(B_j)$,
    combining prior beliefs with observed evidence.
    \item Events $A$ and $B$ are \emph{independent} if
    $\Prob(A \cap B) = \Prob(A)\Prob(B)$; \emph{mutual independence}
    of $n$ events requires the product formula for all $2^n - n - 1$
    subsets of size $\geq 2$.
    \item Pairwise independence does not imply mutual independence:
    the product rule can hold for every pair yet fail for larger subsets.
\end{itemize}

%% ============================================================
%% SECTION 8: FURTHER READING
%% ============================================================

\section{Further Reading}
\label{sec:further_reading}

\begin{itemize}
    \item \textbf{Lesson 10} introduces random variables, probability mass
    functions, and expectation, building directly on the probability space
    framework developed here.
    \item \textbf{Lesson 11} develops common distributions (binomial,
    geometric, Poisson) and their properties using the tools of conditional
    probability and independence.
    \item \textbf{Lesson 12} treats continuous random variables and
    probability density functions, extending the discrete theory.
    \item \textbf{Phase 01, Lesson 1} ($\sigma$-algebras) replaces the
    power-set approach with $\sigma$-algebras, necessary for uncountable
    sample spaces. Phase~01, Lesson~4 revisits probability spaces in the
    measure-theoretic framework.
    \item For a more combinatorial treatment of conditional probability, see
    \citet[Chapter 2]{blitzstein2019}, which includes many worked examples
    from genetics, medicine, and game theory.
    \item The historical development of the Kolmogorov axioms (1933) and
    their unification of discrete and continuous probability is discussed
    in \citet[Chapter 1]{grimmett2020}. Earlier approaches by Laplace,
    Bernoulli, and von Mises lacked the generality of the
    measure-theoretic framework.
\end{itemize}

\bibliography{probability-refs}

\end{document}
